\section{Related Work}
\label{sec:related-work}

\subsection{Quantum computing: foundations and simulators}

The idea that a quantum-mechanical system could simulate physics more efficiently than a
classical computer originates with Feynman~\cite{feynman1982simulating}; Deutsch subsequently
formalized a universal quantum computer~\cite{deutsch1985quantum} and, with Jozsa, gave the
first explicit problem for which a quantum algorithm provably outperforms any classical
deterministic one~\cite{deutschjozsa1992rapid}, the same phase-kickback mechanism
Section~\ref{sec:experiments} demonstrates directly in-world. Shor's
factoring~\cite{shor1997polynomial} and Grover's search~\cite{grover1996fast} algorithms
established quantum computing's practical significance beyond foundational interest.
Nielsen and Chuang's textbook~\cite{nielsenchuang2010} is the standard reference for the
statevector formalism, gate set, and measurement postulate this mod implements directly
(Section~\ref{sec:quantum-simulator}). General-purpose software simulators such as Qiskit and
Cirq, and hardware-backed composers such as IBM's Quantum Composer, let a user assemble a
circuit from a gate palette and inspect the resulting statevector or measurement histogram;
QuantumCraft uses the same underlying representation -- a complex-amplitude vector evolved by
successive unitary matrix multiplication -- but embeds it in a three-dimensional, first-person,
block-based world rather than a schematic canvas, mirroring the relationship CircuitCraft
bears to SPICE-family circuit simulators~\cite{picos2026circuitcraft}.

\subsection{Entanglement and reversible logic}

Two multi-qubit phenomena central to this mod's worked experiments
(Section~\ref{sec:experiments}) have their own foundational literature. Bell's
theorem~\cite{bell1964epr} established that the correlations a pair of entangled qubits
exhibit cannot be reproduced by any local hidden-variable theory, the theoretical result the
Bell-state experiment in Section~\ref{sec:experiments} illustrates operationally rather than
by rederiving the inequality itself; Greenberger, Horne, and Zeilinger's three-qubit
generalization~\cite{ghz1989going} (the ``GHZ state'') sharpens the same point using perfect
rather than merely statistical correlations, and is the target state of this mod's
three-qubit worked experiment. Toffoli's reversible-computing
construction~\cite{toffoli1980reversible} showed that any classical boolean function can be
computed by a reversible (and hence, in particular, quantum-realizable) circuit built from a
single universal three-input gate bearing his name; this mod's Toffoli-gate experiment
(Section~\ref{sec:experiments}) verifies exactly that reversible-AND truth table directly
against a real statevector simulation.

\subsection{Minecraft in education}
\label{sec:minecraft-in-education}

The instructional value of games and simulations has been studied extensively in general: a
meta-analysis by Sitzmann~\cite{sitzmann2011meta} found simulation games to be more effective
for training self-efficacy and knowledge retention than comparison instructional methods
across a range of technical domains. Minecraft occupies a disproportionately large share of
that broader literature. Microsoft maintains and markets a dedicated, classroom-oriented
variant of the game, \emph{Minecraft: Education Edition}, an indication of how extensively the
game had already been adopted informally by educators prior to this commercial
productization. Nebel et al.'s review~\cite{nebel2016mining} surveys this literature broadly
and identifies Minecraft's principal appeal for education as the same property that makes it
appealing as a game: an open, constructive world in which a learner's own artifact, rather
than a worked example, is the object of study, closely matching Papert's constructionist
framing~\cite{papert1980mindstorms}, discussed further in Section~\ref{sec:pedagogy}.

Published case studies span a range of subjects well outside computer science. Bos et
al.~\cite{bos2014mathematics} describe using Minecraft to teach elementary-school mathematics
concepts through building tasks; Short's account~\cite{short2012teaching} covers the use of
the game to teach general science concepts in a secondary-school setting; and Cipollone et
al.~\cite{cipollone2014creative} present a case study of Minecraft used as an open-ended
creative tool. Callaghan's survey~\cite{callaghan2016investigating} identifies a theme common
to each of these studies: teachers report high engagement fairly consistently, whereas
rigorously measured gains in specific learning outcomes are reported less often, a caveat
also present in the general simulation-games literature~\cite{sitzmann2011meta} and one this
paper does not claim to have resolved for QuantumCraft specifically
(Section~\ref{sec:limitations}).

Within that broader literature, the subset closest to the present paper's subject concerns
Minecraft's native redstone subsystem, already a common entry point for teaching classical
digital logic: players build logic gates, adders, and working memory cells entirely from
in-game blocks, and Dezuanni et al.~\cite{dezuanni2015redstone} document children describing
and reasoning about redstone circuits in explicitly electrical language, notwithstanding that
the underlying mechanic is discrete and boolean rather than continuous or quantum-mechanical
(Section~\ref{sec:minecraft-mechanics} discusses redstone's mechanics in more depth). Very
recent work has begun to examine quantum computing specifically as a computer-science
education topic in its own right~\cite{seegerer2021quantum}, independent of any game-based
delivery mechanism. To the author's knowledge, QuantumCraft is the first Minecraft
modification to embed a real statevector quantum-circuit simulator, rather than approximating
quantum behavior with the game's existing discrete, classical redstone mechanics, or
presenting quantum concepts purely through a companion notebook alongside the game.
